\documentclass[9pt,landscape,a4paper]{extarticle}
\usepackage[utf8]{inputenc}
\usepackage[vietnamese]{babel}
\usepackage[T1]{fontenc}
\usepackage{fontspec}
\setmainfont{Libertinus Serif} % As per user's font

% Layout and multicolumn
\usepackage[margin=0.8cm, top=1.5cm, bottom=1.5cm]{geometry}
\usepackage{multicol}
\setlength{\columnsep}{0.8cm}
\setlength{\columnseprule}{0.1pt}

% Math and symbols
\usepackage{amsmath,amssymb}
\usepackage{wasysym}
\usepackage{xcolor}

% Custom Colors (Stanford Cheatsheet style)
\definecolor{stanfordblue}{RGB}{1, 104, 149}
\definecolor{stanfordred}{RGB}{140, 21, 21}
\definecolor{stanfordgreen}{RGB}{0, 100, 0}
\definecolor{cyanbox}{RGB}{220, 245, 255}

% Section styling (mimicking the sheet)
\usepackage{titlesec}
\titleformat{\section}
  {\normalfont\normalsize\bfseries}{\thesection}{1em}{}
\titlespacing*{\section}{0pt}{1.5ex plus 1ex minus .2ex}{0.5ex plus .2ex}

\titleformat{\subsection}
  {\normalfont\small\bfseries}{\thesubsection}{1em}{}
\titlespacing*{\subsection}{0pt}{1ex plus 1ex minus .2ex}{0.5ex plus .2ex}

% Fancy header
\usepackage{fancyhdr}
\pagestyle{fancy}
\fancyhf{}
\renewcommand{\headrulewidth}{0.4pt}
\renewcommand{\footrulewidth}{0.4pt}

% Header content
\lhead{\textsc{MAT3562E} -- \textsc{COMPUTER VISION}}
\rhead{\colorbox{cyan!20}{\makebox[4cm][c]{\textbf{Bùi Quang Chiến -- 23001837}}}}

% Footer content
\lfoot{\textsc{HUS -- VNU}}
\cfoot{\thepage}
\rfoot{\textsc{SPRING 2026}}

% Dense lists
\usepackage{enumitem}
\setlist{nosep, leftmargin=1.2em}
\setlist[itemize,1]{label={\textcolor{stanfordgreen}{\footnotesize$\square$}}}
\setlist[itemize,2]{label=--}

\begin{document}
\begin{center}
    {\Large \bfseries VIP Cheatsheet: LAB07} \\
    \vspace{0.3cm}
    {\large \textsc{Bùi Quang Chiến}} \\
    \vspace{0.2cm}
    {\small April 2026}
\end{center}
\vspace{0.3cm}

\begin{multicols*}{3}

\begin{center}
    {\Large \bfseries VIP Cheatsheet: LAB07} \\
    \vspace{0.3cm}
    {\large \textsc{Bùi Quang Chiến}} \\
    \vspace{0.2cm}
    {\small April 2026}
\end{center}
\vspace{0.3cm}

\begin{multicols*}{3}

\begin{center}
  {\large\bfseries\color{hdr} LAB 07 CHEATSHEET --- Phân Đoạn Biên} \\[1pt]
  {\small Bùi Quang Chiến \textbullet{} 23001837 \textbullet{} MAT3562E Computer Vision 2025--2026}
\end{center}
\hrule\vspace{2pt}

\begin{multicols}{3}

\section{1. Derivative Filters (Đạo Hàm Ảnh Xám)}
\begin{itemize}
    \item \textbf{Cơ bản:} Đo sự chênh lệch $G_x = I(x+1, y) - I(x-1, y)$.
    \item \textbf{Sobel:} \texttt{cv2.Sobel(img, ddepth, dx, dy, ksize)}.\\
        Gradient $\nabla f = [g_x, g_y]^T$.\\
        Biên độ: $M = \sqrt{g_x^2 + g_y^2}$.\\
        Góc: $\theta = \tan^{-1}(g_y/g_x)$.
    \item \textbf{Prewitt / Scharr:} Kernel chập đạo hàm với trọng số khác biệt (chi tiết và nhạy hơn).
\end{itemize}

\section{2. Dò Biên Canny (Canny Edge Detection)}
Thuật toán tối ưu tách nét với 4 bước cốt lõi:
\begin{enumerate}
    \item \textbf{Làm mịn ảnh} (Gaussian Blur) giảm bớt nhiễu xám.
    \item \textbf{Tính toán Gradient} ($M$ và $\theta$).
    \item \textbf{Non-Maximum Suppression:} Khử điểm không cực đại, mỏng hóa biên chỉ giữ 1 nét cực đỉnh sáng liền mạch. 
    \item \textbf{Hysteresis Thresholding:} Lọc ngưỡng kép (\texttt{minVal}, \texttt{maxVal}). Kết nối biên mạnh sang biên yếu, loại bỏ hoàn toàn nhiễu dưới Threshold.
\end{enumerate}
\begin{lstlisting}[language=Python]
edges = cv2.Canny(image, minVal, maxVal)
\end{lstlisting}

\section{3. Chuyển Đổi Hough (Hough Transform)}
Phát hiện hình dạng thông qua không gian tham số.
\begin{tipbox}
\textbf{Đường thẳng (HoughLines):} Tham số $(\rho, \theta)$. 1 đường thẳng = 1 điểm trong không gian Hough. Hàm \texttt{cv2.HoughLinesP} dễ dùng và vẽ đoạn thẳng trực tiếp tốt hơn \texttt{cv2.HoughLines}.
\end{tipbox}
\begin{tipbox}
\textbf{Đường tròn (HoughCircles):} Tham số $(a, b, r)$. Cho phép bắt đường cong, cốc, tiền xu, con ngươi. OpenCV ngầm chạy Canny Edge bên trong với \texttt{cv2.HOUGH\_GRADIENT}.
\end{tipbox}

\section{4. Nối Biên Thành Miền Kín (Edge-To-Region)}
Khi cần tạo object Mask từ Canny Edge:
\begin{lstlisting}[language=Python]
1. edges = cv2.Canny(...)
2. closed = cv2.morphologyEx(edges, cv2.MORPH_CLOSE, kernel) # Boc dut vien
3. contours, _ = cv2.findContours(...)
4. cv2.fillPoly(mask, contours, 255)
\end{lstlisting}

\end{multicols}

\end{multicols*}

\end{multicols*}
\end{document}
